{\scriptsize
\setlength{\tabcolsep}{3pt}
\renewcommand{\arraystretch}{1.5}
\setlist[enumerate]{leftmargin=*, labelsep=0.5em}
\begin{longtable}{|>{\raggedright\arraybackslash}p{0.09\textwidth}|>{\raggedright\arraybackslash}p{0.205\textwidth}|>{\raggedright\arraybackslash}p{0.205\textwidth}|>{\raggedright\arraybackslash}p{0.205\textwidth}|>{\raggedright\arraybackslash}p{0.205\textwidth}|}
\caption{Kế hoạch phân công công việc theo tuần}\\
\hline
\textbf{Giai đoạn} & \textbf{Lưu Chí Lập} & \textbf{Mã Hoàng Linh} & \textbf{Trần Xuân Phú} & \textbf{Nguyễn Hải Đăng} \\
\hline
\endfirsthead

\hline
\textbf{Giai đoạn} & \textbf{Lưu Chí Lập} & \textbf{Mã Hoàng Linh} & \textbf{Trần Xuân Phú} & \textbf{Nguyễn Hải Đăng} \\
\hline
\endhead

\hline
\endfoot

12/1 -- 18/1 & Thiết lập Monorepo cho LingriserSpeaking, cài đặt NestJS Backend và Next.js Frontend. & Thu thập dữ liệu từ các nguồn dataset, tổng hợp 3 tệp JSON với hơn 9.500 mẫu bài viết. & Khởi tạo dự án Next.js, cấu hình TailwindCSS và thiết lập thư viện Shadcn UI. & Khởi tạo Nx Monorepo, thiết lập cấu trúc các microservice cốt lõi: auth, exam, achievement, file-service, gateway. \\
\hline
19/1 -- 25/1 & Hiện thực Database Schema core bằng PostgreSQL, cấu hình Prisma ORM và thực hiện migration đầu tiên. & Tiền xử lý dữ liệu: lọc mẫu không hợp lệ, chuẩn hóa về dạng hội thoại hệ thống -- người dùng -- trợ lý. & Thiết lập Redux Store và quản lý trạng thái cho xác thực, bài thi, người dùng. & Cấu hình thư viện dùng chung: kết nối database, ghi log, xử lý nghiệp vụ theo CQRS, Interceptors, Filters. \\
\hline
26/1 -- 1/2 & Xây dựng Module Auth và Google Auth, tích hợp xác thực đa phương thức. & Cấu hình mô hình Phi-3-mini-4k-instruct, thiết lập QLoRA với 4-bit quantization cho GPU Tesla T4. & Hiện thực giao diện xác thực: biểu mẫu đăng nhập, luồng đăng nhập Email và Google OAuth. & Hiện thực Auth Service: mô hình dữ liệu xác thực, logic đăng ký/đăng nhập, phân quyền theo vai trò. \\
\hline
2/2 -- 8/2 & Hiện thực hệ thống Chat và WebSocket Gateway cho thông báo thời gian thực. & Phân chia tập Train/Eval, bắt đầu Fine-tuning (SFT) trên Kaggle GPU, theo dõi Loss hội tụ. & Xây dựng Layout chính: thanh điều hướng, thanh trên cùng, trang giới thiệu. & Hiện thực API Gateway: cấu hình routing, tích hợp OAuth strategy và cơ chế JWT / Refresh Token. \\
\hline
9/2 -- 22/2 & \multicolumn{4}{c|}{Nghỉ tết Nguyên Đán} \\
\hline
23/2 -- 1/3 & Hiện thực cấu trúc dữ liệu chương trình học: Programs, Courses, Modules. & Hoàn tất Fine-tuning sau 3 Epoch (~7 giờ), Loss đạt mức 0.87. & Hiện thực giao diện luyện thi: màn hình làm bài, thẻ câu hỏi, chọn đề luyện tập. & Hiện thực Exam Service (quản lý): API tạo, sửa, xóa đề thi, phần thi và câu hỏi cho Admin. \\
\hline
2/3 -- 8/3 & Phát triển logic AI-practice và hệ thống Teacher Feedback. & Đánh giá mô hình (MAE, RMSE, Pearson r), so sánh benchmark với gpt-4o-mini. & Hiện thực tính năng Reading và Listening: tô sáng văn bản, trình phát Audio. & Hiện thực Exam Service (thực hành): logic làm bài, lưu đáp án, chấm điểm tự động, lịch sử làm bài. \\
\hline
9/3 -- 15/3 & Xây dựng hệ thống Booking và tích hợp Google Calendar. & Khởi tạo ai-language-evaluator: thiết lập writing server, cấu hình lõi và hệ thống logging. & Hiện thực trang kết quả: giao diện hiển thị điểm số và đáp án chi tiết. & Cấu hình RabbitMQ: gửi/nhận sự kiện bất đồng bộ khi nộp bài và chấm điểm giữa các service. \\
\hline
16/3 -- 22/3 & Hiện thực tính năng phụ huynh theo dõi tiến độ học tập. & Hiện thực logic chấm điểm Writing: Prompt Builder, Regex bóc tách điểm 4 tiêu chí, Guardrails đầu vào. & Hiện thực Dashboard cá nhân: biểu đồ tiến độ học tập, quản lý mục tiêu. & Hiện thực Achievement Service: lắng nghe event RabbitMQ, cấp huy hiệu, tính chuỗi đăng nhập. \\
\hline
23/3 -- 29/3 & Hoàn thiện Backend Dashboard cho Student và Teacher, seed dữ liệu mẫu. & Kết nối RabbitMQ: hiện thực Consumer lắng nghe bài làm và Publisher gửi kết quả chấm điểm. & Hiện thực hệ thống Flashcard và chức năng ghi chú cá nhân. & Cấu hình MinIO/S3, xây dựng API cấp phát presigned-url cho upload/download file. \\
\hline
30/3 -- 5/4 & Kiểm thử và triển khai hệ thống lên Vercel và VPS. & Đẩy mô hình GGUF lên Hugging Face Hub, cấu hình Inference Endpoint với Scale-to-Zero. & Hiện thực Admin Panel: giao diện quản lý đề thi và duyệt đề. & Tối ưu Database: SQL Triggers cho cấu trúc phần thi, Read Models cho truy vấn thống kê. \\
\hline
6/4 -- 12/4 & Triển khai chạy thử tại Đại học Ngoại ngữ Đà Nẵng. & Hiện thực Resource Service: API CRUD cho Blog và Flashcard. & Hiện thực Admin Panel: giao diện quản lý người dùng và bài viết blog. & Hiện thực Report Service: API CRUD báo cáo học tập và thống kê tiến độ. Hỗ trợ merge engapp-v2. \\
\hline
13/4 -- 19/4 & Cải thiện hệ thống theo phản hồi khách hàng. Hỗ trợ merge engapp-v2 vào hệ thống chính. & Hoàn thiện Blog và Flashcard. Tích hợp AI Server với Backend chính, khắc phục lỗi. & Hiện thực Admin Panel: tạo đề thi, cấu trúc phần thi và câu hỏi. & Hiện thực Notification Service và tích hợp Restate workflow cho tác vụ phân tán. \\
\hline
20/4 -- 26/4 & Cải thiện hệ thống theo phản hồi khách hàng. & Kiểm thử toàn trình: Frontend $\rightarrow$ Backend $\rightarrow$ RabbitMQ $\rightarrow$ AI Server $\rightarrow$ UI. Tối ưu thời gian phản hồi. & Tích hợp API Backend, kiểm thử và chỉnh sửa giao diện cho phù hợp với API. & Đóng gói hệ thống bằng Docker. Kiểm thử E2E. \\
\hline
27/4 -- 3/5 & Kiểm thử E2E toàn hệ thống. & Kiểm thử E2E toàn hệ thống. & Kiểm thử E2E toàn hệ thống. & Kiểm thử E2E toàn hệ thống. \\
\hline
4/5 -- 10/5 & Viết báo cáo, chuẩn bị slide bảo vệ. & Viết báo cáo, chuẩn bị slide bảo vệ. & Viết báo cáo, chuẩn bị slide bảo vệ. & Viết báo cáo, chuẩn bị slide bảo vệ. \\
\hline
\end{longtable}
}
